% ============================================================
%  04_clinical_brief.tex
%  Autonomy, Capacity, and the Right to Exit
%  CLINICAL & MEDICAL ETHICS BRIEF
%  For clinical ethics committees, psychiatric professionals,
%  hospital administrators, and medical associations
%  Justin Bogner, 2026
%  Compile: xelatex 04_clinical_brief.tex (twice)
% ============================================================

\documentclass[11pt, letterpaper]{article}
\usepackage{campaign}
\usepackage{xspace}
\usepackage{colortbl}

\campaignheader%
  {Autonomy, Capacity, and the Right to Exit}%
  {Clinical Brief --- Bogner, 2026}

\geometry{
  letterpaper,
  left=1.15in, right=1.15in,
  top=1.1in, bottom=1.2in,
  headheight=22pt
}

\titleformat{\section}
  {\sffamily\large\bfseries\color{C@accent}}
  {}{0em}{}[{\color{C@rulegray}\titlerule[0.5pt]}]

\titleformat{\subsection}
  {\sffamily\normalsize\bfseries\color{C@darkgray}}
  {}{0em}{}

\titleformat{\subsubsection}
  {\sffamily\small\bfseries\itshape\color{C@midgray}}
  {}{0em}{}

\begin{document}

\campaigntitle%
  {AUTONOMY, CAPACITY, AND THE RIGHT TO EXIT}%
  {A Clinical and Medical Ethics Brief on Capacity-Based MAiD}%
  {Clinical / Medical Ethics Brief}%
  {Justin Bogner}%
  {2026}

\suitenote

\begin{policybox}[title={\sffamily\bfseries\small\color{C@accent}
  Purpose and Audience}]
This brief is written for clinical ethics committees, psychiatrists, psychiatric
nurses, hospital administrators, and medical associations engaged with Medical Aid in
Dying policy. It addresses: (1)~the clinical coherence of a capacity-based eligibility
standard; (2)~the assessment protocol proposed under this framework; (3)~the
implications for the therapeutic relationship; (4)~professional liability considerations;
and (5)~specific clinical objections, including the depressive incapacity objection.
\end{policybox}

\bigskip

\section{The Clinical Incoherence of Current Frameworks}

\subsection{The Informed Refusal Doctrine and Its Inconsistent Application}

The foundational principle of modern medical ethics is informed consent: a competent
patient may refuse any treatment, including life-sustaining treatment, on the basis of
their understanding of the treatment and its outcomes, without prior utilization and
without justification to the treating clinician.\footnote{%
  This principle is codified in legislation across jurisdictions (\textit{Health Care
  Consent Act}, Ontario, 1996; \textit{Patient Self-Determination Act}, US, 1990) and
  affirmed in case law. See \textit{Malette v.\ Shulman} (1990), 72 O.R.\ (2d) 417
  (right to refuse emergency treatment); \textit{Re B} [2002] 2 All ER 449 (right to
  refuse ventilator).}

The clinical inconsistency in current MAiD frameworks is precise: this same doctrine is
not applied to the decision to seek self-directed death. A patient may refuse
chemotherapy without trial; they may not, in most jurisdictions, access MAiD without
demonstrating terminal prognosis or externally validated suffering. The asymmetry is
not clinically justified. It reflects historical and legislative contingency, not
principled medical ethics.

\subsection{Suffering as a Clinical Non-Category}

``Unbearable suffering'' is not a clinical concept. It is a legislative construct that
clinical staff are asked to operationalize without any validated instrument or standard
for doing so. The result is that eligibility assessment under suffering-based frameworks
is: (a)~subject to inter-rater unreliability; (b)~sensitive to clinician bias toward or
against MAiD provision; (c)~dependent on the applicant's ability to articulate
subjective distress in terms intelligible to the evaluator.

Capacity, by contrast, is a clinical concept with an established literature, validated
assessment frameworks, and cross-disciplinary consensus (Appelbaum \& Grisso, 1988;
Moye \& Marson, 2007). Replacing suffering with capacity as the eligibility criterion
does not lower clinical standards; it \textit{raises} them by substituting a measurable,
clinically grounded concept for a legally constructed but clinically unmoored one.

\section{The Capacity Assessment Protocol}

\subsection{Operationalizing Capacity for MAiD Assessment}

The Sovereignty Model applies the four-element Appelbaum-Grisso framework to MAiD:

\begin{description}
  \item[Understanding] The applicant can accurately describe what death involves
    (permanent, irreversible cessation), what the proposed provision entails, and what
    the realistic outcomes of all presented alternatives are. Understanding is assessed
    through structured interview, not assumption.

  \item[Appreciation] The applicant understands how the general information applies to
    their specific situation. Appreciation does not require optimism about alternatives;
    it requires demonstrated recognition that the information is relevant to their case.
    The evaluating clinician distinguishes between a patient who says ``SSRIs don't work
    for anyone'' (failure of understanding) and one who says ``SSRIs have a documented
    response rate of approximately 50\% for moderate depression; my history includes two
    adequate trials without response'' (appreciation satisfied).

  \item[Reasoning] The applicant can articulate a coherent account of how their
    understanding of their situation, the alternatives, and their values leads to their
    expressed preference. The reasoning need not be philosophically sophisticated; it
    must be internally coherent and free from gross logical errors or fixed false beliefs.

  \item[Expression of a Consistent Choice] The preference is stable across the
    assessment period and across multiple evaluation contexts. Contextual variation
    (presenting differently to different evaluators) is diagnostically relevant.
    Persistent stability across 90 days and multiple independent evaluations satisfies
    this element.
\end{description}

\subsection{The Evaluating Clinician's Role}

The role of the evaluating clinician under this framework is fundamentally different
from the role assigned under suffering-based frameworks.

\begin{table}[h!]
\centering
\small
\renewcommand{\arraystretch}{1.5}
\begin{tabularx}{\textwidth}{>{\sffamily\bfseries\small}p{3.2cm} X X}
\toprule
\rowcolor{C@accent}
{\color{white}\sffamily\bfseries Dimension} &
{\color{white}\sffamily\bfseries Protectionist Model} &
{\color{white}\sffamily\bfseries Sovereignty Model} \\
\midrule
Primary question
  & Is this person suffering enough?
  & Does this person understand what they are deciding? \\
\rowcolor{C@accentlight}
Role
  & Gatekeeper / adjudicator of suffering
  & Assessor of comprehension and stability \\
Treatment history
  & Required; exhaustion of alternatives typical prerequisite
  & Irrelevant to eligibility; comprehension of alternatives is assessed \\
\rowcolor{C@accentlight}
Clinical instrument
  & None standardized; clinician judgment on suffering
  & Appelbaum-Grisso framework; standardized capacity assessment \\
Liability orientation
  & Institutional (avoid death on our watch)
  & Patient-centered (assess authentic preference) \\
\rowcolor{C@accentlight}
Therapeutic role
  & Custodial retention
  & Restorative alliance \\
\bottomrule
\end{tabularx}
\caption{Comparison of clinician roles under Protectionist and Sovereignty models.}
\end{table}

\subsection{What the Clinical Interview Is Not}

The clinical interview under the Sovereignty Model is not:

\begin{itemize}
  \item A session designed to persuade the patient to choose differently. The clinician
    presents alternatives in sufficient specificity for genuine comprehension; they do
    not advocate for any alternative.
  \item An assessment of whether the clinician agrees with the patient's values or
    finds their reasoning compelling. Clinicians routinely work with patients whose
    values differ from their own without this constituting an eligibility barrier.
  \item An opportunity to require trial of interventions the patient has declined.
    Informed refusal of alternatives is binding. The clinical interview's role is
    to assess whether the refusal is informed, not to override it.
\end{itemize}

\section{The Depressive Incapacity Objection: Clinical Analysis}

This is the objection most frequently raised in clinical contexts, and it deserves
careful clinical analysis rather than brief dismissal.

\subsection{The Objection}

Persons with treatment-resistant depression cannot satisfy the ``appreciation'' criterion
for decisional capacity because hopelessness functions as a cognitive filter: alternatives
are formally understood but experienced as inaccessible. On this account, a capacity
assessment that returns ``capacity present'' in a person with severe depression is a
false positive---the capacity criteria have been formally satisfied, but the authentic
preconditions for autonomous preference are absent.

\subsection{The Clinical Response}

\subsubsection{Depression does not uniformly impair capacity}

The empirical literature does not support the proposition that depression, including
severe depression, uniformly impairs the functional elements of decisional capacity.
Appelbaum and Grisso's foundational studies found that even significantly depressed
inpatients frequently retained capacity for treatment decisions. More recent work has
found that capacity impairment in depression is variable, correlates with severity,
and is not reliably predicted by diagnosis alone.\footnote{%
  See Vollmann et al.\ (2003), ``Competence of mentally ill patients: a comparative
  empirical study,'' \textit{Psychological Medicine}, 33(8): 1463--1471. Also
  Appelbaum \& Grisso (1995), ``The MacArthur Treatment Competence Study.''
  \textit{Law and Human Behavior}, 19(2): 105--126.}

\subsubsection{Appreciation does not require experiential access}

The appreciation criterion requires that the patient understand how the information
applies to their situation. It does not require that they \textit{experience} alternatives
as accessible. A patient who accurately states that SSRIs have a documented response
rate for their condition, that two adequate trials without response are in their record,
and that they decline further trials because the probability of benefit does not, in
their assessment, justify the cost---satisfies the appreciation criterion. The criterion
is epistemic, not phenomenological.

To require, additionally, that the patient \textit{feel hope} that alternatives might
succeed is to import a positive affective requirement that has no basis in the
Appelbaum-Grisso framework and no precedent in clinical capacity doctrine. It also
creates an impossible standard: hopelessness is a symptom of the condition being
assessed; requiring its absence as a precondition for capacity makes the capacity
determination circular.

\subsubsection{The double standard problem}

The depressive incapacity objection, applied consistently, would require that persons
with depression satisfy a higher capacity standard than persons without depression for
all medical decisions. But medicine does not currently impose this: a patient with
severe depression who accepts pharmacological treatment does not have their capacity
for that decision specially scrutinized. The heightened scrutiny applies only when the
decision is to refuse the state-preferred intervention---revealing that the objection
is not about capacity at all, but about the outcome.

\subsection{The Clinical Protocol Response}

The Sovereignty Model does not ignore the depressive incapacity concern. It addresses
it procedurally. A preference that is generated by acute depressive crisis is unlikely
to be stable across 90 days of assessment. A preference that survives the assessment
period---including periods of relative improvement and relative worsening---and
that is articulated consistently across multiple independent evaluations, has survived
the procedural filter that the objection is concerned about. The process is the
response.

\section{Liability Considerations}

\subsection{The Current Liability Model}

Current clinical practice around suicidal ideation is oriented primarily around
institutional liability management: documenting that intervention occurred and that the
clinician met the standard of care for suicide prevention. This orientation produces
the ``liability theater'' documented in the existing literature: coerced safety
contracts of no demonstrated efficacy, pressured hospitalization optimized for
administrative rather than therapeutic outcomes, and discharge into unchanged conditions.

The clinician operating in this model is a liability-reducer first and a therapeutic
ally second. This has consequences for the therapeutic relationship, for patient trust,
and for clinical outcomes.

\subsection{The Sovereignty Model and Liability Reform}

Under the Sovereignty Model, the clinical liability framework shifts. Clinicians who
conduct rigorous, documented capacity assessments following the protocol specified in
this framework---multiple independent evaluations, documented alternatives presentation,
independent coercion screening, 90-day assessment period---are operating within a
structured and documentable standard of care. Liability exposure is reduced by
procedural compliance, not by preventing death at all costs.

The conscientious objection provision is critical here: no clinician is required to
participate in MAiD assessment or provision. The only obligation imposed on objectors
is timely referral to a willing provider. This protects clinical autonomy while
ensuring that patient access is not foreclosed by institutional objection.

\subsection{The Therapeutic Consequence}

The most significant clinical consequence of the Sovereignty Model is the reorientation
of the therapeutic relationship. When the patient possesses an enforceable exit right,
the clinician loses the coercive power of custodial retention. They cannot function as
a jailer mandating biological continuity.

This constraint is not a limitation on clinical effectiveness; it is the precondition
for a different kind of clinical effectiveness. The therapeutic relationship becomes
genuinely collaborative. The clinician's role is to make the patient's life worth
continuing---to address the conditions that are driving the preference, to connect the
patient with resources they may not know exist, to work on the quality of the life
rather than the fact of its continuation. This is a more honest, and ultimately more
effective, clinical orientation.

\section{Assessment Protocol Summary}

\begin{findingbox}[title={\sffamily\bfseries\small\color{C@accent}
  Clinical Protocol at a Glance}]
\small
\textbf{Stage 1 --- Decisional Capacity Assessment} (multiple independent evaluators)
\begin{itemize}
  \item Understanding: permanence of death; realistic outcomes of all alternatives
  \item Appreciation: application of information to own situation
  \item Reasoning: coherent, internally consistent account
  \item Expression: stable preference across assessment period
  \item State assessment: no acute psychosis or crisis-state distortion \textit{currently}
\end{itemize}

\medskip
\textbf{Stage 2 --- Process Comprehension}
\begin{itemize}
  \item Applicant understands why safeguards exist and accepts them
  \item Demonstrates sustained engagement across 90-day minimum period
  \item Failure to engage with Stage 2 is diagnostically relevant to Stage 1
\end{itemize}

\medskip
\textbf{Required Elements of Every Assessment}
\begin{itemize}
  \item Documented alternatives presentation with referrals offered; declinations recorded
  \item Independent coercion screening by unaffiliated social worker or advocate
  \item Financial and dependency circumstances review
  \item Offer of independent legal assistance
  \item Socioeconomic documentation where inadequate support is identified as a driver
\end{itemize}

\medskip
\textbf{Revocation:} Absolute, unconditional, available until the moment of provision.
\end{findingbox}

\section{Clinical Objections: Summary Responses}

\begin{objectionbox}[title={\sffamily\bfseries\small\color{C@warning}
  ``We cannot assess capacity reliably in psychiatric populations''}]
\textbf{Response:} Capacity assessment in psychiatric populations is a well-established
clinical practice with validated instruments (MacArthur Competence Assessment Tool;
Aid to Capacity Evaluation). The proposed framework applies these instruments to a new
context; it does not require new clinical methodology. The challenge of assessing
capacity in complex cases is an argument for rigorous multi-evaluator assessment, not
for categorical exclusion of psychiatric populations from eligibility.
\end{objectionbox}

\medskip

\begin{objectionbox}[title={\sffamily\bfseries\small\color{C@warning}
  ``Psychiatrists will face pressure to approve requests against their clinical judgment''}]
\textbf{Response:} The conscientious objection provision protects clinicians who
decline to participate. The multi-evaluator requirement means no single clinician
carries the assessment alone. The documented capacity protocol creates a professional
standard of care that protects clinicians who conduct rigorous assessments, regardless
of outcome. The pressure concern is addressed by process, not by maintaining a
categorical prohibition.
\end{objectionbox}

\medskip

\begin{objectionbox}[title={\sffamily\bfseries\small\color{C@warning}
  ``This will normalize suicide and undermine suicide prevention''}]
\textbf{Response:} Suicide prevention is directed at crisis-state, impulsive, and
untreated-pathology-driven deaths. The Sovereignty Model is addressed at a categorically
different population: persons with stable, long-considered, informed preferences who
have survived a 90-day assessment process. These populations do not overlap clinically.
Indeed, the explicit, documented nature of MAiD provision---as distinct from unassisted
suicide---may reduce rather than increase the broader social normalization concern,
since it is administered within a visible, regulated, documented clinical process.
\end{objectionbox}

\bigskip
{\small\sffamily\color{C@midgray}
\textbf{Key Clinical References:}
Appelbaum \& Grisso (1988), \textit{NEJM} 319(25);
Appelbaum \& Grisso (1995), MacArthur Treatment Competence Study, \textit{Law and Human Behavior} 19(2);
Beauchamp \& Childress, \textit{Principles of Biomedical Ethics} (8th ed., 2019);
Donnelly \& Campbell, \textit{Journal of Medical Ethics} 50(3), 2024;
Health Canada, \textit{Sixth Annual Report on MAiD} (2025);
Moye \& Marson (2007), \textit{J.\ Gerontology: Psychological Sciences};
Vollmann et al.\ (2003), \textit{Psychological Medicine} 33(8).
}

\end{document}
